\documentclass[12pt]{article}
\usepackage[utf8]{inputenc}
\usepackage[T1]{fontenc}
\usepackage{hyperref}
\title{Wagi Language Resources for K-12 Curriculum Development}
\author{Compiled by Claude Sonnet 5.5 (claude-sonnet-5-5) for Andr\'as Kornai}
\date{}
\begin{document}
\maketitle

\section{Introduction}
Wagi (ISO 639-3: fad; Glottocode wagi1249) is a Papuan language of Madang Province, Papua New Guinea, also known as Kamba, Kauris, Silibob and Mis Furan. Sources give about 3,400 speakers (a 2003 figure); Global Recordings Network gives a population of 1,500, so the figures conflict. Wikitongues rates the language as vigorous (EGIDS 6a). Web search found very little educational material: no primers, textbooks or school programmes were found.

\section{Primary Educational Resources}
No primers, textbooks or school curricula in Wagi were found in search results.

\section{Linguistic and Reference Materials}
\begin{itemize}
\item Wikipedia, Wagi language: general overview. \url{https://en.wikipedia.org/wiki/Wagi_language}
\item Spencer et al.\ 2011, a SIL sociolinguistic survey covering the Kamba, Kauris, Silibob, Foran and Mis dialects, with a lexicon for each, hosted on the Trans New Guinea database. \url{https://transnewguinea.org/source/spencer-et-al-2011?page=3}
\end{itemize}

\section{Religious and Community Resources}
\begin{itemize}
\item Global Recordings Network language page for Wagi (audio recordings of Bible-based material). \url{https://globalrecordings.net/en/language/fad}
\end{itemize}

\section{Digital and Audio Resources}
\begin{itemize}
\item Wikitongues page for Wagi. \url{https://wikitongues.org/languages/fad/}
\item Global Recordings Network, see above. \url{https://globalrecordings.net/en/language/fad}
\end{itemize}

\section{Curriculum Development Context}
The language is reported as vigorous but is not known to be taught in schools. Any curriculum would have to begin from the SIL survey and lexicon data, and would require new materials developed with the community. The conflicting speaker figures should be verified locally.

\end{document}